\documentclass[11pt,a4paper]{article}
\usepackage[utf8]{inputenc}
\usepackage[vietnamese]{babel}
\usepackage{amsmath,amssymb}
\usepackage{geometry}
\geometry{a4paper, total={180mm,260mm}, left=15mm, right=15mm, top=18mm, bottom=18mm}
\usepackage{tikz}
\usetikzlibrary{shapes.geometric, arrows.meta, positioning, calc}
\usepackage{tcolorbox}
\tcbuselibrary{skins,breakable}
\usepackage{xcolor}
\usepackage{multicol}
\usepackage{array}
\usepackage{fontawesome}
\usepackage{fancyhdr}

\pagestyle{fancy}
\fancyhf{}
\renewcommand{\headrulewidth}{0.5pt}
\renewcommand{\footrulewidth}{0.5pt}
\fancyhead[L]{\small\bfseries\color{teal}\faBook\ SỔ TAY KIẾN THỨC TOÁN 3}
\fancyhead[R]{\small\bfseries\color{teal}Lớp Toán Thầy Lĩnh -- loptoanthaylinh.vn}
\fancyfoot[C]{\small\color{gray}-- Sổ tay kiến thức: Bài 53. Luyện tập chung chu vi và diện tích --}

\definecolor{primary}{HTML}{005F73}
\definecolor{secondary}{HTML}{0A9396}
\definecolor{accent}{HTML}{EE9B00}
\definecolor{danger}{HTML}{AE2012}
\definecolor{boxbg}{HTML}{F4F6F8}

\begin{document}

% --- BANNER TIÊU ĐỀ ---
\begin{tcolorbox}[
  enhanced,
  colback=white,
  colframe=primary,
  arc=3mm,
  boxrule=1.2pt,
  left=4mm, right=4mm, top=2.5mm, bottom=2.5mm,
  drop shadow=black!10!white
]
  \centering
  {\Large\bfseries\color{primary}\faBookmark\ BÀI 53: LUYỆN TẬP CHUNG CHU VI VÀ DIỆN TÍCH}\\[1.5mm]
  {\small\itshape\color{black!75}Bảng phân biệt tuyệt đối Chu vi - Diện tích, bài toán ghép hình và ứng dụng thực tế trên 1 trang A4}
\end{tcolorbox}

\vspace{-2mm}


\noindent
\begin{minipage}[t]{0.49\textwidth}\vspace{0pt}
\begin{tcolorbox}[
  enhanced,
  colback=blue!4!white,
  colframe=blue!75!black,
  arc=2mm, boxrule=0.8pt,
  left=2mm, right=2mm, top=1.8mm, bottom=1.8mm,
  title={\footnotesize\bfseries\color{white}\faBalanceScale\ 1. Bảng đối chiếu song song Chu vi -- Diện tích}
]
{\footnotesize
\begin{center}
\renewcommand{\arraystretch}{1.15}
\begin{tabular}{|l|c|c|}
  \hline
  \textbf{Tiêu chí} & \textbf{Hình chữ nhật} & \textbf{Hình vuông} \\
  \hline
  \textbf{Chu vi ($P$)} & $(a + b) \times 2$ & $a \times 4$ \\
  \textbf{Đơn vị $P$} & $\text{cm, m}$ (độ dài) & $\text{cm, m}$ \\
  \hline
  \textbf{Diện tích ($S$)} & $a \times b$ & $a \times a$ \\
  \textbf{Đơn vị $S$} & $\text{cm}^2, \text{m}^2$ (diện tích) & $\text{cm}^2, \text{m}^2$ \\
  \hline
\end{tabular}
\end{center}
\begin{itemize}\setlength{\itemsep}{0.5pt}\setlength{\leftskip}{-2.5mm}
  \item Chu vi là \textbf{độ dài đường viền bao quanh}.
  \item Diện tích là \textbf{toàn bộ bề mặt phẳng bên trong}.
\end{itemize}
}
\end{tcolorbox}
\end{minipage}
\hfill
\begin{minipage}[t]{0.49\textwidth}\vspace{0pt}
\begin{tcolorbox}[
  enhanced,
  colback=teal!4!white,
  colframe=teal!75!black,
  arc=2mm, boxrule=0.8pt,
  left=2mm, right=2mm, top=1.8mm, bottom=1.8mm,
  title={\footnotesize\bfseries\color{white}\faMagic\ 2. Kĩ thuật giải toán ghép và cắt hình}
]
{\footnotesize
\begin{itemize}\setlength{\itemsep}{1.5pt}\setlength{\leftskip}{-2.5mm}
  \item \textbf{Ghép hình:}
  \begin{itemize}\setlength{\itemsep}{0.5pt}\setlength{\leftskip}{-2.5mm}
    \item \textit{Diện tích:} Luôn bằng \textbf{tổng diện tích} các hình nhỏ.
    \item \textit{Chu vi:} \textbf{KHÔNG} bằng tổng chu vi (vì các cạnh dán vào nhau bị ẩn đi bên trong)!
  \end{itemize}
  \item \textbf{Bài toán lát nền nhà / ốp gạch:}\\
  $$\text{Số viên gạch} = \text{Diện tích sàn} : \text{Diện tích 1 viên}$$
\end{itemize}
}
\end{tcolorbox}
\end{minipage}

\vspace{-2.5mm}
% --- KHỐI 3: SƠ ĐỒ TƯ DUY TIKZ ---
\noindent
\begin{tcolorbox}[
  enhanced,
  colback=orange!4!white,
  colframe=orange!80!black,
  arc=2mm, boxrule=0.8pt,
  left=2mm, right=2mm, top=1.5mm, bottom=1.5mm,
  title={\footnotesize\bfseries\color{white}\faSitemap\ SƠ ĐỒ TƯ DUY: BẢN ĐỒ TỔNG KẾT HÌNH HỌC PHẲNG LỚP 3}
]
\centering
\begin{tikzpicture}[
  box/.style={draw, rounded corners=2pt, align=center, font=\scriptsize\bfseries, inner sep=2.5pt, minimum height=0.75cm},
  arr/.style={-{Stealth[length=2mm]}, thick}
]
  \node[box, fill=purple!15, draw=purple!80!black, text width=2.6cm] (h1) at (0, 0) {HÌNH CHỮ NHẬT\\$P = (a + b) \times 2$\\$S = a \times b$};
  \node[box, fill=teal!15, draw=teal!80!black, text width=2.6cm] (h2) at (3.5, 0) {HÌNH VUÔNG\\$P = a \times 4$\\$S = a \times a$};
  \node[box, fill=blue!12, draw=blue!80!black, text width=2.6cm] (h3) at (7.0, 0) {ĐƠN VỊ ĐO\\Chu vi: $\text{cm, m}$\\Diện tích: $\text{cm}^2, \text{m}^2$};
  \node[box, fill=orange!15, draw=orange!80!black, text width=2.6cm] (h4) at (10.5, 0) {ỨNG DỤNG\\Hàng rào $\to$ Chu vi\\Lát gạch $\to$ Diện tích};

  \draw[arr, purple!80!black] (h1) -- (h2);
  \draw[arr, teal!80!black] (h2) -- (h3);
  \draw[arr, blue!80!black] (h3) -- (h4);
\end{tikzpicture}
\end{tcolorbox}

\vspace{-2.5mm}
% --- KHỐI 4: VÍ DỤ & MẸO TRÁNH BẪY ---
\noindent
\begin{minipage}[t]{0.49\textwidth}\vspace{0pt}
\begin{tcolorbox}[
  enhanced,
  colback=green!4!white,
  colframe=green!60!black,
  arc=2mm, boxrule=0.8pt,
  left=2mm, right=2mm, top=1.8mm, bottom=1.8mm,
  title={\footnotesize\bfseries\color{white}\faCheckCircle\ Ví dụ thực hành chuẩn mực}
]
{\footnotesize
\textbf{1. Ghép hình:}\\
Ghép $2$ hình vuông cạnh $4\text{ cm}$ thành một hình chữ nhật.\\
-- Chiều dài: $4 \times 2 = 8\text{ cm}$, chiều rộng: $4\text{ cm}$.\\
-- Chu vi hình chữ nhật: $(8 + 4) \times 2 = \mathbf{24\text{ cm}}$ (nhỏ hơn tổng 2 chu vi là $32\text{ cm}$).\\
-- Diện tích hình chữ nhật: $8 \times 4 = \mathbf{32\text{ cm}^2}$ (bằng tổng 2 diện tích).\\
\textbf{2. Ứng dụng thực tế:}\\
Mảnh đất cần rào xung quanh $\implies$ tính \textbf{Chu vi}.\\
Mặt sàn cần trải thảm $\implies$ tính \textbf{Diện tích}.
}
\end{tcolorbox}
\end{minipage}
\hfill
\begin{minipage}[t]{0.49\textwidth}\vspace{0pt}
\begin{tcolorbox}[
  enhanced,
  colback=red!4!white,
  colframe=red!70!black,
  arc=2mm, boxrule=0.8pt,
  left=2mm, right=2mm, top=1.8mm, bottom=1.8mm,
  title={\footnotesize\bfseries\color{white}\faExclamationTriangle\ Mẹo tránh bẫy trắc nghiệm}
]
{\footnotesize
\begin{itemize}\setlength{\itemsep}{1.5pt}\setlength{\leftskip}{-2.5mm}
  \item \textbf{Bẫy cộng chu vi khi ghép hình:}\\
  Khi ghép các hình lại, các mép giáp nhau không còn là viền ngoài nữa nên chu vi hình mới luôn \textbf{nhỏ hơn} tổng chu vi các hình cũ!
  \item \textbf{Bẫy hỏi một đằng tính một nẻo:}\\
  Đề hỏi "tính chu vi" nhưng lại quen tay áp dụng công thức "diện tích". Hãy đọc kĩ câu hỏi!
\end{itemize}
}
\end{tcolorbox}
\end{minipage}


\end{document}
